\documentclass[11pt]{article}
\usepackage[margin=1in]{geometry}
\usepackage{parskip}
\usepackage{amssymb,amsmath,amsfonts,verbatim,amsthm, amscd}
\usepackage[dvipsnames]{xcolor}
\usepackage[breakable,skins]{tcolorbox}
\usepackage{graphicx,tikz-cd,adjustbox}
\usepackage{float} % breqn can sometimes conflict with amsmath, be mindful
\usepackage{scalerel}
\usepackage{stackengine,wasysym}
\usepackage{lipsum}
\usepackage[backend=biber]{biblatex} % Using biber backend is recommended

\newtcolorbox{mybox}[2][]{
    arc=0mm, enhanced, frame hidden, breakable
}
\usepackage{mathtools}

\usepackage{subcaption}% <-- added

\usepackage{array} % core package
\usepackage{newtxtext,newtxmath} %fuente

\usepackage[cal=boondoxo,scr=euler]{mathalfa}
\usepackage[unicode=true]{hyperref} % permits navigating on the pdf
\hypersetup{pdfencoding=auto}
\usepackage{cleveref} % smart referencing
\usepackage{caption} %

% Define math commands
\newcommand{\calL}{\mathcal{L}}
\newcommand{\calH}{\mathcal{H}}
\newcommand{\bz}{\mathbb{Z}}
\newcommand{\lt}{\normalfont\text{lt}}
\newcommand{\calA}{\mathcal{A}}
\newcommand{\ind}{\text{\normalfont Ind}}
\newcommand{\stab}{\text{\normalfont Stab}}
\newcommand{\hilbM}{\normalfont\text{\underline{H}}_\text{M}}
\newcommand{\aughilbM}{\normalfont\text{H}_\text{M}}
\newcommand{\rank}{\normalfont\text{rk}}
\newcommand{\flats}{\normalfont\text{Flats}}
\newcommand{\flags}{\normalfont\text{Flags}}
\newcommand{\matM}{\normalfont\text{M}}
\newcommand{\hilbMmodF}{\normalfont\text{\underline{H}}_\text{M/F}}
\newcommand{\M}{\normalfont\text{M}}
\newcommand{\N}{\normalfont\text{N}}
\newcommand{\e}{\normalfont\text{e}}
\newcommand{\rk}{\normalfont\text{rk}}
\newcommand{\corank}{\normalfont\text{crk}}
\newcommand{\crk}{\normalfont\text{crk}}

\newcommand{\qbinom}[2]{\begin{bmatrix} #1 \\ #2 \end{bmatrix}_q}

% Theorem environments
\newtheorem{theorem}{Theorem}
\newtheorem{corollary}[theorem]{Corollary}
\newtheorem{lemma}[theorem]{Lemma}
\newtheorem{proposition}[theorem]{Proposition}
\newtheorem{conjecture}[theorem]{Conjecture}
\newtheorem{example}[theorem]{Example}

\theoremstyle{remark}
\newtheorem*{remark}{Remark} % Unnumbered remark
% \theoremstyle{definition} % For proof environment styling - amsthm handles this with \begin{proof}
% \newtheorem*{proof_env}{Proof}

\newcommand{\luis}[1]{{\color{blue} \sf $\clubsuit$ Luis: [#1]}}
\newcommand{\nutan}[1]{{\color{violet} \sf $\clubsuit$ Nutan: [#1]}}
\newcommand{\jacob}[1]{{\color{Bittersweet} \sf $\clubsuit$ Jacob: [#1]}}

\title{\texorpdfstring{Notes (Inverse Kazhdan-Lusztig Polynomials under Deletion)}{Notes (Inverse Kazhdan-Lusztig Polynomials under Deletion)}}
%\author{Nutan Nepal}
%\date{\today}

\begin{document}
\textbf{Notes: Deletion Formula for $Y_M(x)$}\hfill {\small{\today}}

\hrulefill % Horizontal line after title block

Hi everyone,

I'm writing to provide an update on the deletion formula for Y and to ask for some guidance.

Following Tom's suggestion, I've been working on proving the formula by defining the new basis $\mu_F$.
The polynomial $\hat{Y}_M(x)$ is defined as
\[
    \hat{Y}_M(x) = \sum_{F \in \calL(M)} \left( x^{\corank(F)}\mu[F,E] \right) \cdot \hat{Q}_{M^F}(x).
\]



In the initial steps of the proof, I realized that we might need to redefine $\mu_F$ as
\[
    \mu_F = \sum_{G \leq F} x^{\rank(G) - 2\rank(F)}\mu[G,F]\cdot G
\]

so that the polynomial $\hat{Y}_M(x)$ appears in the expression for $\mu_M$ as
\[
    \mu_M = \sum_{G \in \calL(M)} x^{-\rank(G)}\hat{Y}_{M/G}(x^{-2})\cdot \zeta^G
\]




The issue arises when computing $\Delta(\mu_M)$. It seems that my current approach doesn't satisfy the desired property $\Delta(\mu_M) = \mu_{M\setminus e}$.

I have attached a PDF that contains the detailed LaTeX write-up. It outlines the steps I have taken and the calculations.

It's very likely that I am either misunderstanding a subtle part of the original suggestion or overlooking a crucial detail. Tom, I was hoping that by seeing my specific approach, you might be able to identify where I've gone wrong and provide some guidance.

Thank you for your help.

Best regards,

Nutan

\[F\mapsto x^{\rank(F\setminus e)-\rank(F)}\cdot (F\setminus e)\]

\end{document}